\documentclass[11pt]{article}

\usepackage[utf8]{inputenc}
\usepackage[T1]{fontenc}
\usepackage{lmodern}
\usepackage{geometry}
\usepackage{amsmath,amssymb,amsfonts,amsthm,mathtools}
\usepackage{bm}
\usepackage{enumitem}
\usepackage{microtype}
\usepackage{hyperref}
\usepackage{titlesec}
\usepackage{booktabs}
\usepackage{array}
\usepackage{setspace}
\usepackage{mathrsfs}
\usepackage{physics}

\geometry{margin=1in}
\hypersetup{colorlinks=true, linkcolor=black, urlcolor=blue, citecolor=black}
\setlist[enumerate]{topsep=0.4em,itemsep=0.35em}
\setlist[itemize]{topsep=0.3em,itemsep=0.25em}
\titleformat{\section}{\large\bfseries}{\thesection.}{0.5em}{}
\titleformat{\subsection}{\normalsize\bfseries}{\thesubsection.}{0.5em}{}
\setstretch{1.06}

\newcommand{\Aether}{\AE{}ther}
\newcommand{\Aflow}{\AE{}ther-flow}
\newcommand{\TheoryName}{\Aflow{} Interpretation of Relativity}
\newcommand{\TheoryProgram}{\Aether{} / \Aflow{} framework}
\newcommand{\CompletedMetric}{g^{\sharp}}
\newcommand{\CompletionCore}{\mathfrak C^{\mathrm{zr,aff}}}
\newcommand{\TraceCore}{\mathfrak C^{\mathrm{tr}}}
\newcommand{\AffineNucleus}{\mathfrak N^{\mathrm{aff}}}

\newtheorem{definition}{Definition}
\newtheorem{proposition}{Proposition}
\newtheorem{remark}{Remark}
\newtheorem{corollary}{Corollary}
\newtheorem{theorem}{Theorem}

\title{\textbf{The \TheoryName{}}\\[0.4em]
\large Primitive-Reservoir Observer-Reduction\\
\large Minimal Completion-State Shell-Trace Core\\
\large Next Benchmark Criterion}
\author{}
\date{}

\begin{document}

\maketitle

\begin{abstract}
The direct completion-core forcing theorem already removed the exact-shell affine nucleus \(\AffineNucleus\) as the live unexplained stopping point, and the later trace-core collapse theorem has now spent the second admissible move available there by shrinking the full minimal zero-response affine completion core \(\CompletionCore\) to the smaller completion-state shell-trace core \(\TraceCore\). The remaining honest burden is therefore no longer to restate the nucleus, and it is no longer to restate the full completion-core graph presentation.

This note records the routing consequence of that new theorem. The next honest benchmark-facing manuscript must therefore either derive or uniquely force the trace core \(\TraceCore\) from still more primitive explicit substrate structure in a way that changes benchmark-facing status, or prove a genuinely smaller theorem-level collapse, sharper necessity result, or sharper uniqueness result strictly inside \(\TraceCore\). Another same-output deeper-origin relay that merely preserves the same trace core does not count next.

The benchmark boundary remains fixed throughout. The one operative metric is still exactly \(\CompletedMetric\), the ordinary matter route is unchanged, there is no second operative metric, the low-energy matter law is not enlarged, and stronger first-principles promotion language remains paused. The positive result of this note is therefore routing discipline at the new trace-core frontier.
\end{abstract}

\section{Introduction}

The active derivational program of \TheoryName{} is no longer at a stage where the full completion-core graph presentation should be treated as the default unexplained stopping point. The direct completion-core forcing theorem already proved that the exact-shell affine nucleus \(\AffineNucleus\) is forced by the minimal zero-response affine completion core \(\CompletionCore\). The later trace-core collapse theorem then spent the next honest internal move available there: the full affine completion plane and completion body carry no extra benchmark-facing freedom beyond the smaller completion-state shell-trace core \(\TraceCore\).

The present note records the routing consequence of that theorem. It does not reopen the trace-core mathematics, and it does not claim a completed first-principles substrate derivation of gravity. It records only which kind of next manuscript would now count as an honest benchmark-facing gain below the new trace-core frontier.

\paragraph{Framework and claim boundary.}
\TheoryProgram{} denotes the broader \Aether{} / \Aflow{} framework built on the claims that the \Aether{} is the underlying four-dimensional substrate of reality, the \Aflow{} is its intrinsic ordered motion, observed three-dimensional space is the local experiential slice of that deeper substrate, S-time is the experienced order of change arising from matter, light, and the \Aflow{}, and observed expansion is the three-dimensional appearance of deeper four-dimensional ordered motion. Within that framework, \TheoryName{} denotes the exact-closure theory in which the effective gravitational dynamics are adopted to be exactly Einsteinian with universal matter coupling. In the active sequence, \emph{adoption} means use of that established relativistic dynamics without claiming substrate derivation, while \emph{derivation} is reserved for a first-principles recovery from explicit substrate variables.

The present note stays strictly inside that boundary. It records only which kind of next manuscript would now count as an honest benchmark-facing gain below the trace-core frontier.

\section{Fixed Inputs}

\begin{definition}[Trace-core-frontier fixed inputs]
The criterion recorded here uses the following fixed inputs.
\begin{enumerate}
    \item The benchmark exact-closure package, which fixes the public one-metric GR target of the repository.
    \item The benchmark gatekeeping layer, including the post-bridge status correction and the upstream-compression safeguard that blocks same-output deeper-origin relays from counting as new front-line progress once the benchmark-relevant object is unchanged.
    \item The direct exact-shell-affine-nucleus forcing theorem, which already proves that the exact-shell affine nucleus \(\AffineNucleus\) is forced by the completion core \(\CompletionCore\) and therefore removes \(\AffineNucleus\) as an independently unexplained stopping point.
    \item The completion-core trace-collapse theorem, which proves that the full completion-core graph presentation carries no extra benchmark-facing freedom beyond the smaller completion-state shell-trace core \(\TraceCore\).
\end{enumerate}
\end{definition}

\begin{remark}
The present note does not replace those manuscripts. It records the routing consequence of reading them together under the active safeguard against recursive depth drift.
\end{remark}

\section{Why the Trace-Core Collapse Theorem Is the Frontier}

\begin{definition}[Benchmark-neutral trace-core relay]
A \emph{benchmark-neutral trace-core relay} is a manuscript that preserves the same benchmark one-metric output, the same ordinary matter route, the same exact-shell affine nucleus consequence, the same trace core \(\TraceCore\), and the same claim boundary while merely relocating the unexplained substrate layer deeper, restating the full completion-core graph presentation \(\CompletionCore\), restating the nucleus \(\AffineNucleus\), or replaying a preserved deeper ambient-style relay without supplying a new collapse, necessity, uniqueness, or comparably sharp gain inside or below the trace core.
\end{definition}

\begin{proposition}[The trace-core collapse theorem is the live frontier]
At the current stage of the primitive-reservoir observer-reduction line, the theorem collapsing the full completion-core presentation to the smaller completion-state shell-trace core \(\TraceCore\) is the live benchmark-facing frontier below the completion-core forcing theorem.
\end{proposition}

\begin{proof}
That theorem changes benchmark-facing status at current scope because it removes the full completion-core graph presentation as the live internal burden. The current internal bridge datum is no longer the full affine completion-plane package \(\CompletionCore\); it is the smaller trace core \(\TraceCore\). By contrast, a manuscript that preserves the same benchmark package, the same nucleus consequence, and the same trace core while merely pushing the unexplained layer deeper or restating the larger completion-core presentation does not alter benchmark-facing status unless it adds a comparably sharp gain. Therefore the live frontier is the trace-core collapse theorem rather than the earlier completion-core forcing theorem or a later same-output relay.
\end{proof}

\begin{corollary}[Status of the full completion-core presentation as a next-step target]
Under the current routing safeguard, the full minimal zero-response affine completion core \(\CompletionCore\) is no longer itself the default next benchmark-facing target.
\end{corollary}

\begin{proof}
The trace-core collapse theorem already proves that the full completion-core graph presentation is reconstructed from the smaller trace core \(\TraceCore\). Once that is on record, another manuscript that spends the move on the full completion-core presentation again does not remove any further benchmark-relevant freedom. The open burden has therefore moved below that full presentation.
\end{proof}

\section{The Next Benchmark Criterion}

\begin{definition}[Admissible next benchmark-facing trace-core manuscript]
At the current trace-core frontier, a manuscript counts as an admissible next benchmark-facing move only if it does at least one of the following:
\begin{enumerate}
    \item derives or uniquely forces the completion-state shell-trace core \(\TraceCore\) from still more primitive explicit substrate structure in a way that does more than merely relocate the unexplained layer deeper;
    \item proves a genuinely smaller theorem-level collapse, sharper necessity result, or sharper uniqueness result strictly inside \(\TraceCore\);
    \item does either of the previous two while preserving the same benchmark one-metric package, the same ordinary matter route, and the same claim boundary.
\end{enumerate}
\end{definition}

\begin{theorem}[Next honest benchmark-facing task at trace-core scope]
The next honest benchmark-facing manuscript at the current trace-core frontier must either derive or uniquely force the completion-state shell-trace core \(\TraceCore\) from still more primitive explicit substrate structure in a benchmark-relevant way, or else prove a genuinely smaller collapse, sharper necessity result, or sharper uniqueness result inside that trace core. A manuscript that only restates the exact-shell affine nucleus \(\AffineNucleus\), only restates the full completion-core presentation \(\CompletionCore\), or only repeats the same trace-core consequence at a deeper level without such a new gain does not satisfy the criterion.
\end{theorem}

\begin{proof}
By Proposition 1, the live frontier already lies at the theorem that collapses the full completion-core presentation to the smaller trace core. Therefore a second manuscript below it must improve on that status rather than merely accompany it. A deeper-origin manuscript can do so only if it changes benchmark-facing status by more than depth alone. If it simply reproduces the same trace-core consequence through another relay, the routing rule classifies it as benchmark neutral. The remaining honest alternatives are therefore exactly those stated in the definition: either a more primitive derivation or uniqueness forcing of \(\TraceCore\), or a sharper internal collapse, necessity result, or uniqueness result inside \(\TraceCore\).
\end{proof}

\section{What the Next Move Should Not Be}

\begin{proposition}[Screened next moves]
At the current trace-core frontier, none of the following is the default next benchmark-facing move:
\begin{enumerate}
    \item another restatement of the exact-shell affine nucleus \(\AffineNucleus\);
    \item another restatement of the full completion-core graph presentation \(\CompletionCore\);
    \item another same-output deeper-origin relay that merely preserves the same trace-core consequence;
    \item a reopening of preserved deeper ambient-style relays as if they again defined the live burden;
    \item a manuscript that introduces a second operative metric, enlarges the ordinary matter law, or uses stronger promotion language.
\end{enumerate}
\end{proposition}

\begin{proof}
Items (1)--(3) fail the preceding criterion because they do not directly address the current trace-core burden. They revisit either an already-forced bridge object, an already-collapsed larger internal presentation, or the same downstream consequence rather than removing the remaining trace-core freedom or proving a smaller internal datum. Item (4) likewise does not directly address the live burden at current scope. It reopens preserved relays as if they were still frontier-defining rather than settling the present trace-core task. Item (5) violates the fixed claim boundary of the benchmark package and therefore cannot count as a disciplined next step on the active line. The benchmark package remains the exact one-metric GR package already fixed by adoption, so any admissible next manuscript must preserve that boundary.
\end{proof}

\begin{remark}
This screening statement is not a denial that the full completion-core theorem or the preserved deeper ambient-style relays matter historically. It is a routing statement: they are not the honest way to spend the next benchmark-facing move from the current trace-core frontier.
\end{remark}

\section{Conclusion}

The positive result of this note is routing discipline at the current trace-core frontier. The trace-core collapse theorem remains the live benchmark-facing gain because it already removes the full completion-core graph presentation as the internal stopping point on the active line. The full completion-core presentation therefore no longer sets the honest next burden, and neither does the exact-shell affine nucleus that had already been forced one theorem earlier.

The scope remains narrow and honest. The benchmark package is unchanged: the one operative metric is still exactly \(\CompletedMetric\), the ordinary matter route is unchanged, there is no second operative metric, and the low-energy matter law is not enlarged. Nothing here upgrades adoption to derivation or licenses stronger promotion language.

The next open burden is therefore clear. The next benchmark-facing manuscript should derive or uniquely force the completion-state shell-trace core \(\TraceCore\) from still more primitive explicit substrate structure in a way that changes benchmark-facing status, unless the sharper honest gain available instead is a genuinely smaller collapse, sharper necessity result, or sharper uniqueness result strictly inside \(\TraceCore\). That is the clean next manuscript target at current scope.

\end{document}
